% ============================================================
%  Cycle-Accurate Co-Simulation of Computational Delays
%  in Autonomous Vehicle Control
%  Conference submission — double-column, 10pt
% ============================================================
\documentclass[conference]{IEEEtran}

% ---- Packages -----------------------------------------------
\usepackage{amsmath,amssymb,amsthm}
\usepackage{bm}
\usepackage{graphicx}
\usepackage{booktabs}
\usepackage{multirow}
\usepackage{xcolor}
\usepackage{hyperref}
\usepackage{cite}
\usepackage[T1]{fontenc}
\usepackage{tikz}
\usetikzlibrary{arrows.meta,decorations.markings,calc,positioning,patterns}

% ---- Figure colors -------------------------------------------
\definecolor{steelblue}{RGB}{31,119,180}
\definecolor{mutedred}{RGB}{214,39,40}
\definecolor{mutedpurple}{RGB}{148,103,189}
\definecolor{safegreen}{RGB}{44,160,44}
\definecolor{amber}{RGB}{255,195,0}

% ---- Macros --------------------------------------------------
\newcommand{\R}{\mathbb{R}}
\newcommand{\N}{\mathbb{N}}
\newcommand{\RMSE}{\mathrm{RMSE}}
\newcommand{\calX}{\mathcal{X}}
\newcommand{\calU}{\mathcal{U}}
\newcommand{\calH}{\mathcal{H}}
\newcommand{\calW}{\mathcal{W}}
\newcommand{\calT}{\mathcal{T}}
\newcommand{\ZOH}{\mathrm{ZOH}}
\newcommand{\diag}{\mathrm{diag}}
\newcommand{\TODO}[1]{\textcolor{red}{\textbf{[TODO: #1]}}}

% ---- Theorem environments -----------------------------------
\newtheorem{theorem}{Theorem}
\newtheorem{proposition}{Proposition}
\newtheorem{lemma}{Lemma}
\newtheorem{corollary}{Corollary}
\newtheorem{definition}{Definition}
\newtheorem{remark}{Remark}
\newtheorem{assumption}{Assumption}
\newtheorem{example}{Example}

% =============================================================
\begin{document}

\title{SHARC-Drive: Hardware-Aware Simulation for Real-Time Autonomous Driving}

\author{\IEEEauthorblockN{Author~1\IEEEauthorrefmark{1},
    Author~2\IEEEauthorrefmark{2},
    Author~3\IEEEauthorrefmark{1}}
  \IEEEauthorblockA{\IEEEauthorrefmark{1}Affiliation~1 \quad
    \IEEEauthorrefmark{2}Affiliation~2\\
    \texttt{\{email1, email2\}@placeholder.edu}}
}

\maketitle

% =============================================================
\begin{abstract}
  In autonomous driving, control algorithms ranging from PID to
  nonlinear model predictive control (NMPC) must be executed within a
  strict sample period on embedded processors whose computational
  resources vary widely.  The computation time depends on the
  microarchitecture---clock speed, cache size, pipeline depth---and
  interacts with the controller's algorithmic complexity in ways that
  are difficult to predict analytically.  This paper presents a
  co-simulation framework coupling \emph{SHARC} (Simulator for
  Hardware Architecture and Real-time Control) with the \emph{CARLA}
  driving simulator, enabling systematic evaluation of how hardware
  and controller design choices jointly affect closed-loop safety.
\end{abstract}

\begin{IEEEkeywords}
  Hardware--software co-design, autonomous driving, model predictive
  control, computational delay, microarchitecture simulation.
\end{IEEEkeywords}

% =============================================================
\section{Introduction}\label{sec:intro}
An autonomous vehicle's control loop follows a simple pattern:
sense, compute, actuate.  The critical assumption in most control
designs is that computation is instantaneous.  On a fast processor
running a PID controller, this assumption is reasonable---the
computation completes in microseconds.  For a nonlinear model
predictive control (MPC) algorithm on constrained hardware, however,
the computation may approach or exceed the sample period $T$.
When that occurs, the controller misses its deadline: no fresh
command is delivered, and the vehicle continues executing the
previous one.

This coupling between processor capability and control performance
is well understood in
principle~\cite{caccamo2002overruns,findeisen2011robust_delay},
yet the two domains are typically designed in isolation.  Control
engineers assume instantaneous computation; hardware engineers
optimize for generic benchmarks.  Testing on physical prototypes
discovers timing violations late and
expensively~\cite{mihalic2022hil}.

A simulation framework that evaluates the complete loop, plant
physics, control algorithm, \emph{and} processor
microarchitecture, in a single experiment is needed.  This paper
presents such a framework by combining two simulators:
\begin{itemize}
  \item \textbf{SHARC}~\cite{wintz2025sharc} (Simulator for
    Hardware Architecture and Real-time Control) executes a
    controller binary on a parameterized processor model via the
    Scarab cycle-accurate simulator~\cite{scarab}, producing the
    computation time $\tau_k$ at every sample step.
  \item \textbf{CARLA}~\cite{dosovitskiy2017carla}, an open-source
    driving simulator built on Unreal Engine, provides realistic
    vehicle dynamics, urban road networks, and traffic.
\end{itemize}
SHARC is designed around an \emph{abstract dynamics interface}:
any function mapping the current state and control to the next
state can serve as the plant.  We exploit this by replacing the
default ODE integrator with CARLA's physics engine, thereby
obtaining high-fidelity vehicle dynamics without modifying SHARC's
core architecture.\\
\noindent\textbf{Contributions.}
\TODO{}

% =============================================================
\section{Modeling Framework}\label{sec:model}

We develop a general model of a sampled-data system in which
the controller's computation time depends on the state, algorithm,
and hardware.  The model is independent of any particular plant;
CARLA enters in Section~\ref{sec:framework} as one instantiation.

\subsection{Plant and Discretization}

Consider a continuous-time plant
\begin{equation}\label{eq:ct}
  \dot{x} = \tilde{f}(t, x, u, w), \qquad y = h(x),
\end{equation}
with state $x \in \R^{n}$, control input $u \in \R^{m}$, exogenous
input $w \in \R^{p}$, and measured output $y \in \R^{q}$.
A digital controller samples the output at times
$t_k := k\,T$, $k \in \N$, with sample period $T > 0$.
Under zero-order hold, the control is constant over each interval
$[t_k, t_{k+1})$, yielding the discrete-time system
\begin{equation}\label{eq:dt}
  x_{k+1} = f(x_k, u_k, w_k),
\end{equation}
where $f$ encapsulates the continuous dynamics integrated over one
sample period under constant $u_k$ and $w_k := w(t_k)$.

\begin{remark}\label{rem:f_general}
  The map $f$ in~\eqref{eq:dt} need not have a closed-form
  expression.  It may be evaluated by a numerical integrator, a
  lookup table, or, as in our framework, an external physics
  engine.  This generality is central to the design of SHARC and
  enables the CARLA integration described in
  Section~\ref{sec:integration}.
\end{remark}

\subsection{Computation Delay and Zero-Order Hold}

At each sample time $t_k$, the controller reads $y_k = h(x_k)$
and begins computing a new input $\tilde{u}_{k+1}$.  The
computation requires $\tau_k > 0$ seconds.  The computed input is
applied at the \emph{next} sample instant if and only if the
computation finishes within one period; otherwise, the previous
input is held.

\begin{definition}[Zero-Order Hold Under Delay]\label{def:zoh}
  The input applied at sample $k{+}1$ is
  \begin{equation}\label{eq:zoh}
    u_{k+1} = \begin{cases}
      \tilde{u}_{k+1} & \text{if } \tau_k < T, \\
      u_k              & \text{if } \tau_k \geq T.
    \end{cases}
  \end{equation}
\end{definition}

\begin{definition}[Deadline Miss]\label{def:miss}
  At time step $k$, a \emph{deadline miss} occurs if $\tau_k \geq T$.
  The \emph{miss rate} over a horizon of $K$ steps is
  $$\rho := \frac{1}{K}\sum_{k=0}^{K-1}\mathbf{1}[\tau_k \geq T].$$
\end{definition}

\subsection{Hardware-Dependent Computation
  Time}\label{sec:tau_model}

The computation time depends on the algorithm $g$, the current
measurement $y_k$, and the hardware configuration
$\calH \in \R_+^{n_h}$ (clock frequency, cache sizes, pipeline
parameters):
\begin{equation}\label{eq:tau}
  \tau_k = \calT(g, y_k, \calH).
\end{equation}
For a PID controller, $\calT$ is effectively constant: the
computation is a fixed sequence of multiplications and additions.
For MPC, $\calT$ depends on $y_k$ through the optimization
landscape, states far from the reference require more solver
iterations, and on $\calH$ through instruction throughput.

Evaluating $\calT$ analytically is intractable for nontrivial
controllers on realistic processors.  SHARC evaluates it
\emph{empirically} by executing the controller binary inside the
Scarab cycle-accurate simulator for a given hardware configuration.

\subsection{Closed-Loop Dynamics and Cascading Misses}

Substituting~\eqref{eq:zoh} into~\eqref{eq:dt}, the closed-loop
dynamics become
\begin{equation}\label{eq:cl}
  x_{k+1} = f\!\bigl(x_k,\; \mu(x_{0:k},\, \tau_{0:k-1}),\;
  w_k\bigr),
\end{equation}
where $\mu$ is the effective input policy determined by the
ZOH rule.  Because $\tau_k$ depends on $x_k$
via~\eqref{eq:tau}, a feedback loop emerges: the state influences
the computation time, which determines the applied input, which
drives the next state.

\begin{remark}[Cascading deadline misses]\label{rem:cascade}
  For controllers whose computation time grows with the tracking
  error or proximity to constraints, as is typical for iterative optimization-based methods
  such as MPC, deadline misses can cascade.  When a miss occurs,
  the stale input drives the state further from the reference.
  The larger error, in turn, requires more solver iterations at
  the next step, increasing $\tau_{k+1}$ and making a subsequent
  miss more likely.  This positive feedback loop can rapidly
  destabilize the closed loop: a single miss triggers growing
  errors and ever-longer computation times, producing a sharp
  \emph{performance cliff}.  By contrast, controllers with
  state-independent computation time (e.g., PID) are immune to
  this cascade, since $\tau_k$ does not depend on the tracking
  error.
\end{remark}

% =============================================================
\section{Co-Simulation Framework}\label{sec:framework}

\subsection{SHARC}\label{sec:sharc}

SHARC~\cite{wintz2025sharc} is a framework for evaluating
closed-loop control under realistic computational delays.
Its architecture is organized around two user-supplied modules
and one microarchitecture simulator:

\noindent\textbf{Dynamics.}
The user provides a Python class that implements the state
transition $f$ in~\eqref{eq:dt}.  The class exposes two
methods: \texttt{initialize()}, which sets up the plant state,
and \texttt{step(x, u, w)}, which advances the dynamics by
one sample period and returns $x_{k+1}$.  Because the dynamics
are defined in Python, users can implement analytical models
(e.g., an ODE integrator), data-driven surrogates, or, as in
this work, a bridge to an external simulator.

\noindent\textbf{Controller.}
The controller is compiled to a native binary that reads
the measurement $(t_k, y_k, w_k)$ and
generates the control input $\tilde{u}_{k+1}$.
Any control algorithm that can be expressed in C or C++ is
supported by SHARC.

\noindent\textbf{Scarab.}
\TODO{Shengmin: Check and verify this explanation and make it more accurate.}
The Scarab cycle-accurate
simulator~\cite{scarab} executes the controller binary on a user-defined parameterized processor model, producing the computation time
$\tau_k$ in simulated clock cycles.  Scarab models instruction
caches, data caches, branch predictors, out-of-order pipelines,
and memory hierarchy.  The hardware configuration $\calH$ is
specified via a parameter file; sweeping the design space
(e.g., varying clock frequency or cache size) requires only
editing the corresponding parameter.

SHARC provides two execution modes.  In
\emph{execution-driven} (serial) mode, Scarab simulates the
controller instruction-by-instruction during the control loop,
producing $\tau_k$ at each step.  In \emph{trace-driven}
(parallel) mode, the controller runs natively under DynamoRIO
instrumentation~\cite{bruening2003dynamorio} to capture
memory-access traces, which Scarab replays offline; multiple
steps are processed in parallel for faster processing.  When a deadline miss is detected, SHARC re-simulates from the miss point with the corrected state under~\eqref{eq:zoh} as discussed in the SHARC paper \cite{wintz2025sharc}.

\subsection{CARLA Integration}\label{sec:integration}

CARLA~\cite{dosovitskiy2017carla} is an open-source driving
simulator with realistic vehicle dynamics, urban road networks,
and traffic participants.  We integrate CARLA into SHARC by
implementing a Python dynamics class that connects to
CARLA's synchronous API.  At initialization, the class
spawns the ego vehicle and NPC traffic, enables synchronous
mode with physics step $\Delta t = T$, and seeds the random
number generators for reproducibility.  At each step, the
class applies the control input to CARLA, ticks the world,
and reads back the new state.  Because $\Delta t = T$, each
CARLA tick corresponds to exactly one discrete-time step
in~\eqref{eq:dt}.  From SHARC's perspective, CARLA is simply
a black-box instantiation of the abstract~$f$. \TODO{Yash \& Malavikha: explain our determinism framework, explain why we simulate from the beginning for each batch.}

% =============================================================
\section{Controller Formulations}\label{sec:controllers}

We evaluate a few controllers with qualitatively different
computational profiles.

\subsection{State and Control Representation}\label{sec:state}

The ego vehicle state and control input are
\begin{align}
  x_k &:= (p_x,\, p_y,\, \psi,\, v)^\top
    \in \R^4, \label{eq:state}\\
  u_k &:= (\alpha,\, \delta,\, \beta)^\top
    \in [0,1] \times [-1,1] \times [0,1],\label{eq:control}\\
    w_k &:= (w_x^+, w_y^+)^\top \in \R^2 \label{eq:exo}
\end{align}
where $(p_x, p_y)$ is position, $\psi$ is heading (yaw), $v$ is
speed, $(w_x^+, w_y^+)$ is the next waypoint provided by CARLA from the HD map, and
$(\alpha, \delta, \beta)$ are throttle, steering, and brake.
The exogenous input $w_k := (w_x^+, w_y^+)^\top$ provides the
reference waypoint to the controller at each step.

\subsection{PID Controller}\label{sec:pid}

The PID controller decomposes vehicle control into two independent
SISO loops.

\noindent\textbf{Longitudinal.}  Track a constant reference speed
$v^*$:
\begin{equation}\label{eq:pid_lon}
  \sigma_k = K_P^\ell\, e_k^v + K_I^\ell \sum_{j=0}^{k} e_j^v\, T
    + K_D^\ell\, \frac{e_k^v - e_{k-1}^v}{T},
\end{equation}
where $e_k^v := v^* - v_k$.  Throttle and brake are
$\alpha_k = \mathrm{clip}(\sigma_k,\, 0,\, 1)$,\;
$\beta_k = \mathrm{clip}(-\sigma_k,\, 0,\, 1)$.

\noindent\textbf{Lateral.}  Steer toward the current waypoint.
The heading error is
$e_k^\psi := \mathrm{atan2}(w_y^+ - p_y,\; w_x^+ - p_x) - \psi_k$,
and the steering command is
\begin{equation}\label{eq:pid_lat}
  \delta_k = \mathrm{clip}\!\bigl(K_P^r\, e_k^\psi + K_I^r
    \textstyle\sum_j e_j^\psi\, T + K_D^r\,
    (e_k^\psi - e_{k-1}^\psi)/T,\;\pm 1\bigr).
\end{equation}
PID requires a fixed number of arithmetic operations per step,
so $\tau_k \approx \mathrm{const}$.

% ---- TikZ figure: MPC waypoint tracking --------------------
\begin{figure}[t]
  \centering
  \resizebox{0.45\textwidth}{!}{%
  \begin{tikzpicture}[
      >=Stealth,
      % ---- Style definitions ----
      wp/.style={circle, fill=steelblue, draw=steelblue!60!black,
        line width=0.5pt, inner sep=2.2pt},
      pred/.style={circle, fill=mutedred, draw=mutedred!60!black,
        line width=0.5pt, inner sep=2.6pt},
      car/.style={rectangle, draw=black!85, line width=0.9pt,
        fill=gray!90, minimum width=15mm, minimum height=7.5mm,
        rounded corners=2pt},
      dist/.style={densely dashed, mutedpurple, line width=0.75pt},
      refpath/.style={steelblue, line width=1.6pt, opacity=0.40},
      predpath/.style={mutedred, line width=1.2pt, densely dotted},
      veloc/.style={-{Stealth[length=5pt]}, line width=1.3pt, safegreen!80!black},
      refline/.style={black!35, dashed, line width=0.6pt},
      anglearc/.style={steelblue!65!black, line width=0.8pt,
        -{Stealth[length=4pt]}},
      annot/.style={font=\small},
    ]

    % ================================================================
    % Reference waypoints along an S-curve
    % ================================================================
    \foreach \i/\wx/\wy in {
       % 1/0.0/0.00,   2/0.5/0.19,   3/1.0/0.44,   4/1.5/0.74,
       5/2.0/1.08,   6/2.5/1.42,   7/3.0/1.74,   8/3.5/2.02,
       9/4.0/2.24,  10/4.5/2.40,  11/5.0/2.48,  12/5.5/2.50,
      13/6.0/2.46,  14/6.5/2.36,  15/7.0/2.22,  16/7.5/2.04,
      17/8.0/1.84,  18/8.5/1.62,  19/9.0/1.40} {
      \node[wp] (w\i) at (\wx,\wy) {};
    }

    % ---- Waypoint labels (first, middle, last) ----
    \node[annot, steelblue!80!black, above=3pt] at (w5)  {$w^{(1)}$};
    \node[annot, steelblue!80!black, above=3pt] at (w10) {$w^{(i)}$};
    \node[annot, steelblue!80!black, below=3pt] at (w19) {$w^{(N_w)}$};

    % ---- Smooth reference path ----
    \draw[refpath]
      plot[smooth, tension=0.50] coordinates {
        (0.0,0.00) (0.5,0.19) (1.0,0.44) (1.5,0.74)
        (2.0,1.08) (2.5,1.42) (3.0,1.74) (3.5,2.02)
        (4.0,2.24) (4.5,2.40) (5.0,2.48) (5.5,2.50)
        (6.0,2.46) (6.5,2.36) (7.0,2.22) (7.5,2.04)
        (8.0,1.84) (8.5,1.62) (9.0,1.40)
      };

    % ================================================================
    % Ego vehicle — heading angle 35 deg, position (0.80, 0.22)
    % ================================================================
    \node[car, rotate=20] (ego) at (0.80, 0.22) {};

    % Label: state vector below the vehicle (absolute coords)
    % \node[annot, black!70, anchor=north] at (0.80, -0.28)
    %   {\textbf{Ego}:\; $x_k = (p_x,\,p_y,\,\psi,\,v)^\top$};

    % ---- East-direction reference dashed line (for angle arc) ----
    \draw[refline] (ego.center) -- ++(0:1.55);

    % ---- Heading angle arc: 0 deg -> 35 deg, radius 1.1 ----
    % Arc starts at (ego.center + (1.1, 0)) and sweeps to 35 deg
    \draw[anglearc]
      ($(ego.center)+(1.10,0)$) arc (0:20:1.10)
      node[midway, below right=-8pt and 1pt, annot,
        steelblue!70!black] {$\psi_k$};

    % ---- Velocity vector (direction of motion, separate from arc) ----
    % \draw[veloc] (ego.center) -- ++(35:1.55)
    %   node[above right=1pt and 0pt, annot, safegreen!70!black] {$v_k$};

    % ================================================================
    % MPC predicted states — offset BELOW reference path so
    % distance lines d_{j|k} are clearly non-zero
    % ================================================================
    \node[pred] (p1) at (2.20, 0.75) {}; % near w5=(2.0,1.08)
    \node[pred] (p2) at (3.35, 1.35) {}; % near w7=(3.0,1.74)
    \node[pred] (p3) at (4.45, 1.86) {}; % near w9=(4.0,2.24)
    \node[pred] (p4) at (5.55, 2.12) {}; % near w11=(5.0,2.48)
    \node[pred] (p5) at (6.65, 2.05) {}; % near w13=(6.0,2.46)

    % ---- Predicted trajectory ----
    \draw[predpath] (ego.center) -- (p1) -- (p2) -- (p3) -- (p4) -- (p5);

    % ---- Predicted state labels ----
    \node[annot, mutedred, right=2pt of p1]  {$\hat{x}_{k+1|k}$};
    \node[annot, mutedred, below=0pt of p3]  {$\hat{x}_{k+3|k}$};
    \node[annot, mutedred, below=0pt of p5]  {$\hat{x}_{k+N_p|k}$};

    % ================================================================
    % Distance lines to nearest waypoints
    % ================================================================
    \draw[dist] (p1) -- (w5)
      node[midway, above right=-2pt, annot, mutedpurple] {$d_{1|k}$};
    \draw[dist] (p2) -- (w7)
      % node[midway, left=3pt, annot, mutedpurple] {$d_{2|k}$};
    \draw[dist] (p3) -- (w10)
      node[midway, right=0pt, annot, mutedpurple] {$d_{3|k}$};
    \draw[dist] (p4) -- (w12)
      % node[midway, left=3pt, annot, mutedpurple] {$d_{4|k}$};
    \draw[dist] (p5) -- (w14)
      node[midway, above=3pt, annot, mutedpurple] {$d_{5|k}$};

    % ================================================================
    % Legend box (positioned below and to the right of path)
    % ================================================================
    \begin{scope}[shift={(6.65,-0.75)}]
      \draw[black!25, rounded corners=4pt, fill=white, fill opacity=0.92]
        (-0.25,-0.40) rectangle (3.55, 1.20);

      % Row 1: waypoint
      \node[wp] at (0.10, 0.95) {};
      \node[font=\scriptsize, anchor=west] at (0.30, 0.95)
        {reference waypoint $w^{(i)}$};

      % Row 2: predicted state
      \node[pred] at (0.10, 0.60) {};
      \node[font=\scriptsize, anchor=west] at (0.30, 0.60)
        {predicted state $\hat{x}_{k+j|k}$};

      % Row 3: predicted trajectory
      \draw[predpath] (0.0, 0.25) -- (0.25, 0.25);
      \node[font=\scriptsize, anchor=west] at (0.30, 0.25)
        {predicted trajectory};

      % Row 4: distance to reference
      \draw[dist] (0.0,-0.10) -- (0.25,-0.10);
      \node[font=\scriptsize, anchor=west] at (0.30,-0.10)
        {tracking distance $d_{j|k}$};

      % % Row 5: velocity vector
      % \draw[veloc] (0.0,-0.45) -- (0.25,-0.45);
      % \node[font=\scriptsize, anchor=west] at (0.30,-0.45)
      %   {velocity $v_k$};

      % % Row 6: heading angle
      % \draw[anglearc] (0.18,-0.72) arc (0:45:0.18);
      % \node[font=\scriptsize, anchor=west] at (0.30,-0.72)
      %   {heading angle $\psi_k$};
    \end{scope}

  \end{tikzpicture}}
  \caption{MPC waypoint-tracking formulation.  The ego vehicle
    at position $(p_x, p_y)$ moves with speed $v_k$ and heading
    $\psi_k$ (measured from east, shown as an arc).
    Reference waypoints $\{w^{(i)}\}_{i=1}^{N_w}$ (\textcolor{steelblue}{blue})
    define the desired path.  At each step $k$, MPC rolls out the
    kinematic bicycle model~\eqref{eq:bicycle} to obtain
    $N_p$ predicted states $\hat{x}_{k+j|k}$ (\textcolor{mutedred}{red})
    and minimizes the sum of squared distances $d_{j|k}$
    (\textcolor{mutedpurple}{purple dashed}) from each predicted
    position to the closest waypoint~\eqref{eq:closest_wp}.}
  \label{fig:mpc_waypoints}
\end{figure}

\subsection{Model Predictive Controller}\label{sec:mpc}
\TODO{Tyler: Check and verify this part and make it more accurate.}
The MPC controller receives a vector of $N_w$ reference waypoints
$\mathcal{W} := \{w^{(i)}\}_{i=1}^{N_w}$,
$w^{(i)} = (w_x^{(i)}, w_y^{(i)})$,
from a route planner and solves a finite-horizon optimal control
problem at each sample step.  The formulation is illustrated in
Fig.~\ref{fig:mpc_waypoints}.

\noindent\textbf{Prediction model.}  We use a kinematic bicycle
model for internal prediction:
\begin{equation}\label{eq:bicycle}
  \hat{x}_{j+1|k} = \hat{x}_{j|k}
  + T \begin{bmatrix}
    v_{j|k} \cos\psi_{j|k} \\
    v_{j|k} \sin\psi_{j|k} \\
    (v_{j|k} / L)\tan\delta_{j|k} \\
    a_{j|k}
  \end{bmatrix},
\end{equation}
where $L$ is the wheelbase and $a_{j|k}$ is the longitudinal
acceleration derived from throttle and brake.

\noindent\textbf{Cost function.}  The cost penalizes the
Euclidean distance from each predicted position to the
\emph{closest} waypoint, along with speed tracking and control
effort:
\begin{equation}\label{eq:mpc_cost}
  J = \sum_{j=1}^{N_p} \bigl[
    q_\ell\, d_{j|k}^2
    + q_v\, (v_{j|k} - v^*)^2
  \bigr]
  + \sum_{j=0}^{N_c-1} \|u_{j|k}\|_R^2,
\end{equation}
where
\begin{equation}\label{eq:closest_wp}
  d_{j|k} := \min_{i \in \{1,\ldots,N_w\}}
  \bigl\|(p_{x,j|k},\, p_{y,j|k}) - w^{(i)}\bigr\|
\end{equation}
is the distance to the closest waypoint, $v^*$ is the reference
speed, $q_\ell, q_v > 0$ are weights, and $R \succ 0$ penalizes
control effort.

\noindent\textbf{Optimization problem.}  At each $t_k$:
\begin{subequations}\label{eq:mpc_opt}
\begin{align}
  \min_{u_{0|k},\ldots,u_{N_c-1|k}} \quad & J \\
  \text{s.t.} \quad
    & \hat{x}_{j+1|k} = F(\hat{x}_{j|k}, u_{j|k}), \\
    & \hat{x}_{0|k} = x_k, \label{eq:mpc_ic}\\
    & u_{j|k} \in \calU, \qquad j = 0,\ldots,N_c{-}1,
\end{align}
\end{subequations}
with prediction horizon $N_p$, control horizon $N_c \leq N_p$,
and input constraint set $\calU$.  Only the first input
$u_{0|k}$ is applied (receding horizon);
the result becomes $\tilde{u}_{k+1}$
in~\eqref{eq:zoh}.

\begin{remark}[State-dependent computation time]\label{rem:mpc_tau}
  The number of solver iterations, and hence $\tau_k$, depends
  on the current state.  Near the reference, the optimization is
  well-conditioned; after a miss, the perturbed state requires more
  iterations.  This is the mechanism underlying the cascading
  misses described in Remark~\ref{rem:cascade}.
\end{remark}

\begin{remark}[Computational complexity ordering]
  PID requires $O(1)$ arithmetic operations per step, whereas
  NLMPC iteratively solves a nonconvex NLP whose size grows
  with the prediction horizon $N_p$.  The ordering
  $\tau_k^{\mathrm{PID}} \ll \tau_k^{\mathrm{NLMPC}}$
  is central to our hardware sensitivity analysis.
\end{remark}

% =============================================================
\section{Experimental Setup}\label{sec:experiments}

\subsection{Driving Scenarios}

Table~\ref{tab:scenarios} summarizes the scenarios drawn from
CARLA's built-in maps and town configurations.  Scenarios are ordered
by \emph{safety criticality}: the higher the row, the less margin is
available for delayed actuation. \TODO{Fill this table correctly}

\begin{table}[h]
  \centering
  \caption{Driving scenarios used for evaluation.}
  \label{tab:scenarios}
  \renewcommand{\arraystretch}{1.2}
  \footnotesize
  \begin{tabular}{llp{2.6cm}c}
    \toprule
    \textbf{Scenario} & \textbf{Map} & \textbf{Key challenge}
      & \textbf{$T$ (s)}\\
    \midrule
    Straight-road cruising     & Town04 & Speed tracking        & 0.25 \\
    Highway lane following     & Town04 & Lateral precision     & 0.25 \\
    Adaptive cruise control    & Town04 & Lead-car gap          & 0.20 \\
    Intersection crossing      & Town03 & Timing, collision av. & 0.15 \\
    Emergency braking          & Town01 & Min. stopping dist.   & 0.10 \\
    High-speed race lap        & Town04 & Stability at speed    & 0.10 \\
    \bottomrule
  \end{tabular}
\end{table}

\subsection{Hardware Configurations}\label{sec:hw_configs}

We evaluate each controller on a set of processor configurations
by varying the clock frequency and L1 data cache size while
holding all other microarchitectural parameters (pipeline width,
branch predictor, memory latency) at a fixed baseline.
Table~\ref{tab:hw} summarizes the configurations.

\begin{table}[t]
  \centering
  \caption{Processor configurations evaluated in all experiments.
    Each row defines one Scarab parameter set.}
  \label{tab:hw}
  \footnotesize
  \begin{tabular}{clcl}
    \toprule
    \textbf{Config} & $f_{\text{clk}}$ (MHz)
      & $S_{\text{L1}}$ (kB) & \textbf{Description} \\
    \midrule
    A & 2000 &  64  & Fast baseline  \\
    B & 1000 &  64  & Halved frequency  \\
    C & 500  &  64  & Quarter frequency  \\
    D & 250  &  64  & Eighth frequency  \\
    E & 2000 &   8  & Small cache  \\
    F & 2000 & 128  & Large cache  \\
    G & 500  &   8  & Constrained  \\
    \bottomrule
  \end{tabular}
\end{table}

\subsection{Performance Metrics}

Let $\{x_k\}_{k=0}^K$ be the simulated trajectory and
$\{\bar{x}_k\}_{k=0}^K$ the reference.

\noindent\textbf{Position RMSE:}
\begin{equation}\label{eq:rmse}
  \RMSE_{\mathrm{pos}} :=
  \sqrt{\frac{1}{K}\sum_{k=0}^{K-1}
    \bigl[(p_x^k - \bar{p}_x^k)^2 + (p_y^k - \bar{p}_y^k)^2\bigr]}.
\end{equation}

\noindent\textbf{Speed RMSE:}
$\RMSE_v := \sqrt{\frac{1}{K}\sum_{k=0}^{K-1}(v_k - v^*)^2}$.

\noindent\textbf{Deadline miss rate} $\rho$ (Definition~\ref{def:miss}).

\noindent\textbf{Mean computation delay:}
$\bar{\tau} := \frac{1}{K}\sum_{k=0}^{K-1} \tau_k$, with distribution
$p(\tau_k)$ visualized as a box plot.

\noindent\textbf{Safety events:} number of (a) collisions with the
environment, (b) lane-boundary violations ($|e_y| > 1.5$\,m), and
(c) speed-limit exceedances ($v > 1.2\,v^*$) per episode.
% =============================================================
\section{Results}\label{sec:results}

We report results for lane following with $T = 0.1$\,s and
$K = 64$ steps.  Table~\ref{tab:results} presents the key
metrics for each hardware configuration.

\begin{table}[t]
  \centering
  \caption{Experimental results: mean computation delay
    $\bar{\tau}$, deadline miss rate $\rho$, position RMSE,
    and speed RMSE for each controller and hardware
    configuration ($T = 0.1$\,s, $K = 64$\,steps).} 
  \label{tab:results}
  \footnotesize
  \renewcommand{\arraystretch}{1.15}
  \begin{tabular}{c l rr rr}
    \toprule
    & & \multicolumn{2}{c}{\textbf{PID}}
      & \multicolumn{2}{c}{\textbf{MPC}} \\
    \cmidrule(lr){3-4}\cmidrule(lr){5-6}
    \textbf{Cfg}
      & \textbf{Metric}
      & \multicolumn{1}{c}{Value}
      & \multicolumn{1}{c}{$\rho$}
      & \multicolumn{1}{c}{Value}
      & \multicolumn{1}{c}{$\rho$} \\
    \midrule
    \multirow{2}{*}{A}
      & $\bar{\tau}$ ($\mu$s)         & 14   & \multirow{2}{*}{0.0}
                                      & ---  & \multirow{2}{*}{---} \\
      & $\RMSE_{\text{pos}}$ (m)      & ---  &
                                      & ---  &  \\
    \midrule
    \multirow{2}{*}{B}
      & $\bar{\tau}$ ($\mu$s)         & 28   & \multirow{2}{*}{0.0}
                                      & ---  & \multirow{2}{*}{---} \\
      & $\RMSE_{\text{pos}}$ (m)      & ---  &
                                      & ---  &  \\
    \midrule
    \multirow{2}{*}{C}
      & $\bar{\tau}$ ($\mu$s)         & 56   & \multirow{2}{*}{0.0}
                                      & ---  & \multirow{2}{*}{---} \\
      & $\RMSE_{\text{pos}}$ (m)      & ---  &
                                      & ---  &  \\
    \midrule
    \multirow{2}{*}{D}
      & $\bar{\tau}$ ($\mu$s)         & 112  & \multirow{2}{*}{0.0}
                                      & ---  & \multirow{2}{*}{---} \\
      & $\RMSE_{\text{pos}}$ (m)      & ---  &
                                      & ---  &  \\
    \midrule
    \multirow{2}{*}{G}
      & $\bar{\tau}$ ($\mu$s)         & 56   & \multirow{2}{*}{0.0}
                                      & ---  & \multirow{2}{*}{---} \\
      & $\RMSE_{\text{pos}}$ (m)      & ---  &
                                      & ---  &  \\
    \bottomrule
  \end{tabular}
\end{table}

\subsection{Computation Delay Distributions}

\TODO{Figure: Box plots of $\tau_k$ distributions for PID, MPC
across different hardware and scenarios }

\subsection{Tracking Performance vs.\ Hardware}
\TODO{Figure: RMSE$_\text{pos}$ vs.\ clock frequency sweep for
NLMPC and LMPC on the intersection and emergency-braking scenarios.
Expected finding: step-function degradation at a critical clock—the
``hardware--performance cliff.''  PID flat across all frequencies,
confirming computational simplicity as robustness.}

\TODO{Figure: Pareto frontier plots: RMSE$_\text{pos}$ vs.\ clock
speed for LMPC and NLMPC with varying prediction horizon $N_p \in
\{5, 10, 20\}$.  Each $(N_p, f_\text{clk})$ pair is a design point.
Expected finding: smaller $N_p$ reduces computation time, lowers
$\rho$, but increases tracking error; the Pareto frontier identifies
optimal $(N_p, f_\text{clk})$ tradeoffs for each AV scenario.}

This analysis operationalizes the co-design objective: rather than
choosing the largest $N_p$ that a given processor can \emph{nominally}
run, the SHARC--CARLA framework reveals the largest $N_p$ the processor
can run \emph{within the sample period} of a specific driving scenario,
accounting for worst-case instruction-level timing.



\subsection{Cache Sensitivity}
\TODO{Cache sensitivity plots}

% =============================================================
\section{Discussion}\label{sec:discussion}

\TODO{This is just a placeholder for now.}
\noindent\textbf{The performance cliff.}
The most significant finding in SHARC-based
experiments~\cite{wintz2025sharc} is that control degradation is
not gradual.  A small reduction in processor capability can flip
the system from zero misses to cascading misses
(Remark~\ref{rem:cascade}).  Our CARLA integration extends this
observation to a realistic driving scenario with high-fidelity
vehicle dynamics and traffic.
\noindent\textbf{MPC horizon as a co-design variable.}
The classical MPC design question ``how large should $N_p$ be?'' is
answered in closed-loop terms by our Pareto analysis.  On constrained
hardware, halving $N_p$ can eliminate all deadline misses while
incurring only a modest RMSE increase—a trade that the designer cannot evaluate without SHARC.

\noindent\textbf{Limitations.}
CARLA's physics engine, while detailed, is not a certified vehicle
model. The current study evaluates only two controllers; extending to a broader family of algorithms is future work.

% =============================================================
\section{Conclusion}\label{sec:conclusion}
\TODO{}
% =============================================================
\bibliographystyle{IEEEtran}
\bibliography{references}

\end{document}
